% =========================
% Guidance / Requirements (section header comment)
% =========================
\noindent\textbf{What this section should contain.}
Specify the empirical method and experimental protocol. Define baselines and comparison targets, and explain which comparisons are feasible.
If some comparisons are not feasible, provide a clear rationale and narrow claims accordingly.

\begin{itemize}
  \item \textit{[ACM]} \emph{Empirical methodology options:} action research (researchers intervene in a real organization using the artifact);
  case study (researchers observe a real organization using the artifact); controlled experiment (human participants use the artifact);
  quantitative simulation (artifact assessed, often against a competing artifact, in an artificial environment);
  benchmarking study (artifact assessed using one or more benchmarks); or another method with a clear and convincing rationale.

  \item \textit{[ACM]} Clearly indicate which of the above empirical methodologies is used.

  \item \textit{[ACM]} Compare the approach to a justified and appropriate baseline.

  \item \textit{[ACM]} \emph{Comparative evaluation requirement:} either empirically compare the artifact to one or more state-of-the-art alternatives,
  or compare it to one or more state-of-the-art benchmarks, or provide a clear and convincing rationale for why comparative evaluation is impractical.

  \item \textit{[NeurIPS]} Main experimental results include the new method and baselines.

  \item \textit{[NeurIPS]} If only a subset of experiments are reproducible, explicitly state which ones are reproducible.

  \item \textit{[Project-specific]} \emph{Add any additional local requirements here (e.g., internal checklists, ablation expectations, artifact release constraints).}
\end{itemize}
